\section{From Task-Specific Prediction to Reusable Representation}

Chapters 1 and 2 introduced learning as function approximation and deep learning as the practice of learning representations rather than hand-crafting features.
Chapter 3 then emphasized that architecture matters because it imposes inductive bias.
This chapter studies a further modern shift:
we increasingly train models not only to solve a single task, but to acquire internal structure that can be reused across many tasks.

\medskip

The central thesis of the chapter is simple but important.
Modern machine learning increasingly learns \emph{reusable internal coordinates}.
A good representation organizes raw inputs into a space in which many downstream tasks become easier to express, learn, and adapt.
This is the conceptual core of pretraining, transfer learning, self-supervision, and foundation models.

\subsection*{From a single predictor to a representation-plus-head decomposition}

In the most classical supervised-learning picture, one trains a predictor
\[
f : X \to Y
\]
directly for one specified task.
Given training examples \((x_i,y_i)\), the objective is to choose a function from a hypothesis class that minimizes empirical risk and hopefully generalizes.
This perspective remains correct, but it hides an important internal factorization.
Many modern systems are better viewed as
\[
f = g \circ \phi,
\]
where
\[
\phi : X \to Z
\]
is a representation map and
\[
g : Z \to Y
\]
is a downstream predictor.

\begin{definition}[Representation]
Let \(X\) be an input space and \(Z\) a feature or latent space.
A \emph{representation} is a map
\[
\phi : X \to Z
\]
that assigns to each input \(x \in X\) a coded description \(z = \phi(x)\).
The representation is useful for a task if a predictor acting on \(z\) can solve that task more simply, more robustly, or with less labeled data than a predictor acting directly on raw inputs.
\end{definition}

This definition is deliberately broad.
The representation might be a hidden layer of a neural network, an embedding vector, a sequence of token features, a set of patch embeddings, or a learned state summarizing a time series.
The key point is that the representation is not itself the final answer.
It is an intermediate coordinate system in which answers become easier to extract.

\medskip

Once we write \(f = g \circ \phi\), the learning problem splits into two conceptual parts:
\begin{itemize}
    \item \(\phi\) learns how to organize the input space,
    \item \(g\) learns how to read out the information needed for a particular task.
\end{itemize}
This decomposition is one of the deepest structural ideas in modern machine learning.
It says that part of what is learned should be \emph{reusable}.

\subsection*{Task-specific features versus transferable features}

A feature is \emph{task-specific} if it is useful mainly for one narrowly defined prediction problem.
For example, if one trains a classifier only to distinguish cats from dogs, the internal representation may organize the data in a way highly tuned to that distinction.
That can work well for the original task while being of limited use for retrieval, captioning, segmentation, or medical anomaly detection.

By contrast, a feature is \emph{transferable} if it participates in many tasks.
Edges, shapes, object parts, syntax, semantic relations, speaker characteristics, and latent dynamical state are all examples of structure that can be useful beyond one label space.
A representation is therefore more transferable when it captures regularities of the data-generating process rather than merely memorizing the decision boundary of one benchmark.

A useful formal way to express this is through a family of tasks.
Let \(\mathcal{T}\) index tasks, with each task \(t \in \mathcal{T}\) having target space \(Y_t\), distribution \(P_t\) over \(X \times Y_t\), and downstream predictor class \(\mathcal{G}_t\) acting on \(Z\).
The same representation \(\phi\) is used for all tasks, while the readout \(g_t\) may vary with the task.

\begin{definition}[Transferable representation relative to a task family]
Let \(\{(P_t,\mathcal{G}_t): t \in \mathcal{T}\}\) be a family of downstream learning problems.
A representation \(\phi : X \to Z\) is \emph{transferable} relative to this family if, for many tasks \(t\), there exists a comparatively simple predictor \(g_t \in \mathcal{G}_t\) such that
\[
R_t(g_t \circ \phi)
\]
is small, where \(R_t\) denotes the population risk on task \(t\).
\end{definition}

This definition is intentionally qualitative in one place: what counts as ``simple'' depends on context.
Sometimes simplicity means linear probes.
Sometimes it means small sample complexity, shallow heads, or minimal fine-tuning.
The main idea is stable: a good representation turns many tasks into easy readout problems.

\subsection*{A fixed representation induces a new hypothesis class}

Once \(\phi\) is fixed, the downstream learner no longer searches over arbitrary functions on \(X\).
Instead it searches over functions on \(Z\), then composes them with \(\phi\).
This induces a precise new hypothesis class.

\begin{proposition}[A fixed representation induces a hypothesis class over \(X\)]
Let \(\phi : X \to Z\) be fixed, and let \(\mathcal{G}\) be a class of predictors \(g : Z \to Y\).
Define
\[
\mathcal{H}_{\phi}
=
\{g \circ \phi : g \in \mathcal{G}\}.
\]
Then \(\mathcal{H}_{\phi}\) is a hypothesis class on \(X\).
Moreover, empirical risk minimization over \(\mathcal{G}\) on transformed examples \((\phi(x_i),y_i)\) is exactly empirical risk minimization over \(\mathcal{H}_{\phi}\) on the original examples \((x_i,y_i)\).
\end{proposition}

\paragraph{Proof sketch.}
For every \(g \in \mathcal{G}\), the composition \(g \circ \phi\) is a function from \(X\) to \(Y\), so \(\mathcal{H}_{\phi}\) is a valid hypothesis class on \(X\).
Now consider the empirical objective
\[
\frac{1}{n}\sum_{i=1}^n \ell(g(\phi(x_i)),y_i).
\]
This is numerically identical to
\[
\frac{1}{n}\sum_{i=1}^n \ell(h(x_i),y_i)
\quad\text{with}\quad h = g \circ \phi.
\]
Thus minimizing over \(g \in \mathcal{G}\) on transformed data is the same optimization problem as minimizing over \(h \in \mathcal{H}_{\phi}\) on the original data.
\qed

This proposition may look obvious, but it is conceptually powerful.
A representation is not just a convenient hidden layer.
It changes the effective hypothesis class available to the learner.
When \(\phi\) is well chosen, difficult nonlinear structure in \(X\) may become easy linear or low-complexity structure in \(Z\).
When \(\phi\) is poorly chosen, even simple downstream tasks may remain hard.

\begin{example}[Linear probes as a restricted readout class]
If \(\mathcal{G}\) is the class of affine maps on \(Z = \mathbb{R}^m\), then
\[
\mathcal{H}_{\phi}
=
\{x \mapsto w^{\top}\phi(x) + b : w \in \mathbb{R}^m,\; b \in \mathbb{R}\}.
\]
This is the hypothesis class induced by \emph{linear probing}.
The probe is linear in representation space even if the resulting predictor is highly nonlinear as a function of the original input \(x\).
\end{example}

So linear probing is not a contradiction.
It is simple in \(Z\), but it can correspond to a sophisticated decision rule in \(X\) because the nonlinearity has been delegated to the representation map.

\subsection*{Why reusable representations are valuable}

The task-specific paradigm asks us to collect labeled data and solve each task separately.
That strategy becomes expensive when tasks are numerous or labels are costly.
In many realistic domains, however, unlabeled or weakly labeled data are abundant.
Language is everywhere.
Images and video are easy to collect.
Sensor streams, scientific measurements, and interaction traces can be massive.
What is scarce is often not data itself but task-specific supervision.

This asymmetry motivates pretraining.
One first learns a representation on broad data, often using a generic objective such as prediction, masking, contrast, reconstruction, or next-token modeling.
One then adapts the model to a downstream task by training a small head, fine-tuning part of the model, or fine-tuning all of it.
The reusable component is the representation.

\begin{figure}[h]
\centering
\fbox{%
\begin{minipage}{0.93\textwidth}
\centering
\textbf{Pretrain--then--adapt pipeline}

\vspace{0.6em}
\begin{tabular}{c}
Large raw dataset \(\{x\}\) \\
\(\downarrow\) pretraining objective \\
Learn shared representation map \(\phi_{\theta} : X \to Z\) \\
\(\downarrow\) reuse across tasks \\
\begin{tabular}{ccc}
Task A head \(g_A\) & Task B head \(g_B\) & Task C head \(g_C\)
\end{tabular} \\
\(\downarrow\) \\
Predictors \(g_A \circ \phi_{\theta}\), \(g_B \circ \phi_{\theta}\), \(g_C \circ \phi_{\theta}\)
\end{tabular}
\end{minipage}}
\caption{A schematic picture of the modern pipeline.
The shared representation is learned first and then adapted or probed for multiple downstream tasks.
The central claim of representation learning is that the same internal coordinates can support many different predictors.}
\end{figure}

The pipeline can be realized in different ways.
Sometimes the representation is frozen and only a probe is trained.
Sometimes the whole model is fine-tuned.
Sometimes adaptation is done by prompts, low-rank updates, or small inserted modules.
The engineering details differ, but the conceptual structure is the same:
learn reusable internal organization first, specialize later.

\subsection*{A worked example: linear probing after representation learning}

A toy example makes the geometry explicit.
Suppose our input space is \(X = \mathbb{R}^2\), and the classification problem is to distinguish points on an inner circle of radius \(1\) from points on an outer circle of radius \(2\).
In the original space \(X\), these classes are not linearly separable by a straight line.
No affine function
\[
x \mapsto a^{\top}x + b
\]
can separate one circle from the other.

Now define the representation
\[
\phi(x) = \|x\|^2.
\]
This maps each point in \(\mathbb{R}^2\) to a scalar in \(Z = \mathbb{R}\).
For points on the inner circle,
\[
\phi(x) = 1,
\]
and for points on the outer circle,
\[
\phi(x) = 4.
\]
A linear probe on \(Z\) takes the form
\[
g(z) = wz + b.
\]
Choose \(w = 1\) and \(b = -\tfrac52\).
Then
\[
g(1) = -\tfrac32 < 0,
\qquad
 g(4) = \tfrac32 > 0.
\]
Thus the classifier
\[
\hat y = \mathrm{sign}(g(\phi(x)))
\]
perfectly separates the two classes.

\begin{example}[What the probe is really telling us]
The downstream classifier is linear only in the representation space \(Z\), not in the original input space \(X\).
The nonlinear geometric work was done by \(\phi\), which collapsed radial information into a one-dimensional coordinate.
This is why linear probes are informative.
If a simple probe succeeds, it suggests that the representation has already organized the relevant task information into an accessible form.
\end{example}

This example is intentionally small, but it captures a general pattern.
A representation can convert a hard decision boundary in raw input space into a simple separator in latent space.
Modern deep learning uses very high-dimensional and learned versions of the same idea.
Instead of hand-defining \(\phi(x)=\|x\|^2\), the model learns a representation from data.
A linear or shallow downstream head then tests what information has become explicit.

\subsection*{What is modern about this picture}

It is tempting to think that representation learning is merely an implementation detail inside neural networks.
That would miss the historical shift.
Earlier machine learning often treated features as external to the learner: one engineered them first and then trained a predictor.
Deep learning already changed this by learning internal features jointly with the task.
The more recent development goes further still:
large-scale pretraining seeks features that are not tied to one task at all.
The representation becomes a reusable computational resource.

This is why the language of ``foundation models'' emerged.
A foundation model is not defined only by parameter count or benchmark accuracy.
Its defining ambition is to provide a broadly reusable representational substrate from which many downstream capabilities can be adapted.
That ambition raises scientific questions as well as engineering ones.
What kinds of pretraining objectives induce transferable structure?
How should one measure transferability?
When does a representation support linear probing, and when does it require deeper adaptation?
What parts of a learned representation are robust, spurious, or task-dependent?
These are central themes of the rest of the chapter.

\paragraph{What to remember.}
A representation is a map \(\phi : X \to Z\) that organizes raw data into a space where downstream prediction becomes easier.
Modern predictors are often best understood as compositions \(g \circ \phi\).
Once \(\phi\) is fixed, the downstream learner operates in an induced hypothesis class \(\mathcal{H}_{\phi} = \{g \circ \phi : g \in \mathcal{G}\}\).
Transferable representations are those that allow many tasks to be solved by relatively simple readouts.
The pretrain--then--adapt paradigm is built on the claim that reusable internal coordinates can be learned once and reused many times.

\paragraph{Common confusion.}
A reusable representation is not the same thing as a universally optimal representation.
Transfer is always relative to a family of tasks and a notion of downstream simplicity.
It is also easy to think that linear probing means the original problem was linear.
That is false:
a problem can be highly nonlinear in \(X\) while becoming linearly separable after a learned map into \(Z\).
Finally, pretraining does not guarantee transfer.
A representation can be excellent for one distribution or objective and poor for another.

\paragraph{Exercises.}
\begin{enumerate}
    \item Let \(\phi : X \to Z\) be fixed and let \(\mathcal{G}\) be a downstream hypothesis class on \(Z\).
    Prove carefully that empirical risk minimization over \(\mathcal{G}\) on transformed data is equivalent to empirical risk minimization over \(\mathcal{H}_{\phi} = \{g \circ \phi : g \in \mathcal{G}\}\) on the original data.
    \item Give an example of a feature that is useful for one classification task but unlikely to transfer well to a broader family of tasks.
    Explain why it is task-specific rather than reusable.
    \item In the concentric-circle example, verify that no affine classifier on \(\mathbb{R}^2\) can separate the inner circle from the outer circle, but that a threshold on \(\|x\|^2\) can.
    \item Suppose \(\phi : \mathbb{R}^d \to \mathbb{R}^m\) is learned and a downstream task uses a linear probe.
    Write the resulting predictor explicitly and explain why it may still define a nonlinear decision boundary in the original input space.
    \item Compare two adaptation strategies after pretraining: freezing \(\phi\) and training only \(g\), versus fine-tuning both \(\phi\) and \(g\).
    What are the likely advantages and risks of each?
\end{enumerate}

\paragraph{Notes and references.}
For broad perspectives on representation learning, see Bengio, Courville, and Vincent on representation learning and Goodfellow, Bengio, and Courville on deep learning.
For transfer learning as a general paradigm, Pan and Yang remain a standard reference.
The diagnostic use of linear classifiers on top of learned features appears throughout the modern literature on probing and representation evaluation.
Historically, the conceptual shift in this section leads directly to self-supervised learning, large-scale pretraining, and foundation-model thinking:
the common thread is the attempt to learn internal coordinates that outlive any one benchmark task.
